% SPDX-License-Identifier: PMPL-1.0-or-later
% arcvix-octad-data-model.tex — arXiv-style academic paper on the Octad Data Model
% VeriSimDB: Drift-Aware Multi-Modal Data Identity in Federated Systems
%
% Author: Jonathan D.A. Jewell
% Affiliation: Independent Researcher / Hyperpolymath
% Date: March 2026

\documentclass[11pt,twocolumn]{article}

% ── Packages ──────────────────────────────────────────────────────────────────
\usepackage[utf8]{inputenc}
\usepackage[T1]{fontenc}
\usepackage{amsmath,amssymb,amsthm}
\usepackage{mathtools}
\usepackage{stmaryrd}          % ⟦ ⟧ semantic brackets
\usepackage{booktabs}
\usepackage{tabularx}
\usepackage{graphicx}
\usepackage{hyperref}
\usepackage{cleveref}
\usepackage{algorithm}
\usepackage{algpseudocode}
\usepackage{listings}
\usepackage{xcolor}
\usepackage{tikz}
\usetikzlibrary{arrows.meta,positioning,shapes.geometric,calc,fit}
\usepackage[margin=1in]{geometry}
\usepackage{microtype}
\usepackage{enumitem}
\usepackage{caption}
\usepackage{subcaption}
\usepackage{fancyhdr}
\usepackage{xspace}

% ── Theorem environments ──────────────────────────────────────────────────────
\newtheorem{definition}{Definition}[section]
\newtheorem{theorem}{Theorem}[section]
\newtheorem{lemma}[theorem]{Lemma}
\newtheorem{proposition}[theorem]{Proposition}
\newtheorem{corollary}[theorem]{Corollary}
\newtheorem{invariant}{Invariant}[section]
\newtheorem{property}{Property}[section]

% ── Macros ────────────────────────────────────────────────────────────────────
\newcommand{\verisimdb}{\textsc{VeriSimDB}\xspace}
\newcommand{\octad}{\mathcal{O}}
\newcommand{\modset}{\mathcal{M}}
\newcommand{\driftfn}{\Delta}
\newcommand{\identity}{\mathcal{I}}
\newcommand{\repairfn}{\mathcal{R}}
\newcommand{\policyfn}{\mathcal{P}}
\newcommand{\queryfn}{\mathcal{Q}}
\newcommand{\sembrack}[1]{\llbracket #1 \rrbracket}

% ── Listing style ─────────────────────────────────────────────────────────────
\definecolor{codebg}{gray}{0.96}
\definecolor{codegreen}{rgb}{0.0,0.5,0.0}
\definecolor{codegray}{rgb}{0.5,0.5,0.5}
\definecolor{codeblue}{rgb}{0.1,0.2,0.6}
\lstset{
  basicstyle=\ttfamily\footnotesize,
  backgroundcolor=\color{codebg},
  keywordstyle=\color{codeblue}\bfseries,
  commentstyle=\color{codegreen},
  stringstyle=\color{codegray},
  numbers=left,
  numberstyle=\tiny\color{codegray},
  breaklines=true,
  frame=single,
  framesep=3pt,
  xleftmargin=1.5em,
  framexleftmargin=1.5em,
  captionpos=b
}

% ── Title ─────────────────────────────────────────────────────────────────────

\title{%
  \textbf{The Octad Data Model: Drift-Aware Multi-Modal\\
  Data Identity in Federated Systems}%
}

\author{
  Jonathan D.A.\ Jewell\\
  \textit{Independent Researcher}\\
  \texttt{j.d.a.jewell@open.ac.uk}
}

\date{March 2026}

% ══════════════════════════════════════════════════════════════════════════════
\begin{document}
\maketitle

% ── Abstract ──────────────────────────────────────────────────────────────────
\begin{abstract}
Modern data systems face a fundamental tension between the richness of
multi-modal representations and the coherence guarantees required by federated
environments.  An entity may simultaneously be a node in a knowledge graph, a
dense vector embedding, a multi-dimensional tensor, a semantically typed
artefact, a full-text document, a versioned temporal object, a provenance-tracked
lineage chain, and a spatially located datum.  Existing multi-model databases
typically store these facets in loosely coupled subsystems with no unified
identity or consistency semantics.

We present the \emph{Octad Data Model}, the core abstraction of \verisimdb, in
which every entity exists simultaneously across exactly eight \emph{modalities}
--- Graph, Vector, Tensor, Semantic, Document, Temporal, Provenance, and Spatial
--- bound to a single UUID-based identity.  Rather than enforcing strict
consistency through conflict-free replicated data types or blockchain-style
consensus, \verisimdb treats cross-modal \emph{drift} as a first-class concept:
drift is detected through statistical and formal methods, classified by type and
severity, and resolved through configurable policy-driven repair strategies that
range from fully automatic to human-supervised.  We formalise the octad algebra,
define drift metrics across modality pairs, and describe a federated
verification protocol suitable for sovereign data stores that cannot be centrally
controlled.  We report on a Rust-based implementation using Oxigraph, Tantivy,
HNSW indexing, and R-tree spatial indexing, and present experimental results on
drift detection accuracy and cross-modal query performance over synthetic and
real-world datasets.
\end{abstract}

\noindent\textbf{Keywords:}
multi-model databases, data identity, drift detection, federated systems,
knowledge graphs, vector search, temporal databases, provenance tracking

% ── 1. Introduction ──────────────────────────────────────────────────────────
\section{Introduction}\label{sec:intro}

The contemporary data landscape is characterised by an explosion of
heterogeneous representations.  A single scientific publication may be
represented as an RDF triple in a bibliographic knowledge graph, a 768-dimensional
dense embedding in a semantic search index, a tensor of citation features in a
recommendation model, a full-text document in a search engine, a versioned
artefact in a temporal database, and a geolocated entity in a spatial index.
Each representation captures a distinct \emph{modality} of the same underlying
entity, yet in practice these modalities are stored in separate, independently
managed systems with no shared notion of identity or consistency.

This fragmentation creates three classes of problems:

\begin{enumerate}[leftmargin=*]
  \item \textbf{Identity fragmentation.}  The same entity acquires different
    identifiers in each system, requiring brittle reconciliation logic.
  \item \textbf{Semantic drift.}  When one modality is updated (e.g., a
    paper is retracted), other modalities may continue to serve stale or
    contradictory information indefinitely.
  \item \textbf{Federation barriers.}  Sovereign institutions cannot share
    data across organisational boundaries without either ceding control to a
    central authority or accepting the chaos of uncoordinated peer-to-peer
    replication.
\end{enumerate}

Existing multi-model databases --- ArangoDB~\cite{arangodb2023}, SurrealDB~\cite{surrealdb2024},
TerminusDB~\cite{terminusdb2023} --- address identity fragmentation to varying
degrees but treat consistency as a binary property: either strongly consistent
(via MVCC, Raft, or CRDTs) or eventually consistent with no formal drift
semantics.  Decentralised systems such as IPFS~\cite{benet2014ipfs} and
Solid~\cite{sambra2016solid} provide content-addressable storage and
access-control mechanisms, but offer no cross-modal coherence guarantees.

We argue that the missing abstraction is \emph{drift-aware multi-modal
identity}: a data model in which (a) an entity's identity is preserved across
all representational modalities, (b) divergence between modalities is detected,
classified, and quantified continuously, and (c) resolution of detected drift is
governed by explicit, auditable policies rather than implicit convergence
assumptions.

This paper makes the following contributions:

\begin{enumerate}[leftmargin=*]
  \item We formalise the \emph{Octad Data Model}, in which each entity exists
    simultaneously across eight modalities, bound by a single identity and
    subject to algebraic invariants (\cref{sec:octad}).
  \item We define a \emph{drift calculus} with formal metrics for detecting
    and classifying divergence across modality pairs (\cref{sec:drift}).
  \item We describe a \emph{policy-driven repair framework} that supports
    automatic, semi-automatic, and manual resolution strategies with immutable
    audit trails (\cref{sec:repair}).
  \item We present a \emph{cross-modal query} language that permits type-safe
    access across modality boundaries (\cref{sec:query}).
  \item We report on a production-quality implementation in Rust and Elixir,
    together with experimental evaluation on drift detection accuracy and
    query latency (\cref{sec:implementation,sec:evaluation}).
\end{enumerate}

The remainder of this paper is organised as follows.  \Cref{sec:background}
surveys relevant background.  \Cref{sec:octad} presents the formal octad model.
\Cref{sec:drift} defines drift detection and classification.  \Cref{sec:repair}
describes policy-driven repair.  \Cref{sec:query} introduces cross-modal
querying.  \Cref{sec:implementation} details the implementation.
\Cref{sec:evaluation} reports experimental results.  \Cref{sec:related} discusses
related work, and \cref{sec:conclusion} concludes.

% ── 2. Background ────────────────────────────────────────────────────────────
\section{Background}\label{sec:background}

\subsection{Multi-Model Databases}

The term \emph{multi-model database} denotes a system that natively supports
more than one data model (e.g., document, graph, key-value) within a single
engine~\cite{lu2019multimodel}.  ArangoDB~\cite{arangodb2023} combines
document, graph, and key-value access patterns over a shared storage engine.
SurrealDB~\cite{surrealdb2024} extends this with record links, full-text
search, and vector similarity.  OrientDB provides document-graph duality.
These systems achieve multi-model access by embedding multiple query
interpreters atop a common storage layer, but they do not formalise the
relationship \emph{between} models for the same entity.  An entity's graph
representation and its document representation are conceptually independent;
no mechanism detects when they diverge.

\subsection{Graph Databases and Knowledge Graphs}

RDF-based systems (e.g., Apache Jena, Oxigraph~\cite{oxigraph2023},
Blazegraph) and property-graph systems (e.g., Neo4j, JanusGraph) represent
entities as nodes connected by typed edges.  SPARQL and Cypher provide
declarative query interfaces.  While graph representations excel at capturing
relational structure, they lack native support for dense vector retrieval,
temporal versioning, or spatial indexing, requiring external systems for these
modalities.

\subsection{Vector Stores and Dense Retrieval}

The rise of transformer-based embeddings has driven the development of
purpose-built vector databases: Pinecone, Weaviate, Milvus, Qdrant, and
Chroma.  These systems index high-dimensional vectors using approximate
nearest-neighbour algorithms such as HNSW~\cite{malkov2018hnsw} or
IVF~\cite{jegou2011pq}.  While some offer metadata filtering, they provide
no formal mechanism for detecting when an embedding becomes stale relative
to the source document it was derived from.

\subsection{Temporal Databases and Version Control}

Temporal databases~\cite{snodgrass1999temporal} maintain the history of data
changes, supporting time-travel queries (``what was the value at time $t$?'')
and bitemporal models (transaction time vs.\ valid time).  TerminusDB~\cite{terminusdb2023}
applies version-control semantics (branches, merges, diffs) to graph data.
Dolt provides git-style versioning for relational tables.  However, temporal
guarantees in these systems apply within a single modality; cross-modal
temporal coherence (e.g., ensuring a vector embedding corresponds to the
correct document version) is left to application logic.

\subsection{Provenance and Lineage Tracking}

The W3C PROV model~\cite{w3cprov2013} defines a vocabulary for describing the
provenance of data entities, activities, and agents.  Systems such as
Apache Atlas and Amundsen provide metadata lineage for data warehouses.
Content-addressable storage (IPFS~\cite{benet2014ipfs}, Git) ensures
immutability through cryptographic hashing.  These approaches track
\emph{where data came from} but do not address \emph{whether} concurrent
modality representations remain consistent with the provenance record.

\subsection{Spatial Databases}

PostGIS, SpatiaLite, and H3~\cite{uber2018h3} provide geospatial indexing
using R-trees, quadtrees, or hexagonal hierarchical grids.  The OGC Simple
Features standard and WGS84 coordinate reference system are widely adopted.
Spatial databases rarely integrate with graph or vector retrieval, creating
yet another silo.

\subsection{FAIR Principles and Federated Data}

The FAIR principles~\cite{wilkinson2016fair} (Findable, Accessible,
Interoperable, Reusable) have become the standard for open-science data
management.  Solid~\cite{sambra2016solid} provides user-controlled data pods
with linked-data access control.  The European Open Science Cloud (EOSC)
and GAIA-X pursue federated data spaces.  These initiatives address access
control and discoverability but lack formal models for cross-modal identity
and drift.

% ── 3. The Octad Data Model ──────────────────────────────────────────────────
\section{The Octad Data Model}\label{sec:octad}

\subsection{Design Rationale}

The central observation motivating the octad model is that multi-modal data
representations are not independent: they are \emph{projections} of a single
underlying entity onto different representational spaces.  A knowledge graph
edge stating ``paper $p$ cites paper $q$'' and a vector embedding of $p$ that
places it near $q$ in embedding space are not independent facts --- they are
correlated views that should remain mutually consistent.  The octad model
makes this correlation explicit and enforceable.

We identify eight modalities as the minimal complete basis for contemporary
data-intensive systems.  This choice is informed by a survey of production
workloads spanning scientific publishing, geospatial intelligence, financial
compliance, and AI/ML pipelines.  While any finite enumeration is necessarily
incomplete, we demonstrate that these eight modalities cover the representational
needs of each surveyed domain.

\subsection{Formal Definition}

\begin{definition}[Modality Set]\label{def:modset}
The \emph{modality set} $\modset$ is the fixed eight-element set:
\[
  \modset = \{
    \mathsf{G}, \mathsf{V}, \mathsf{T}, \mathsf{S},
    \mathsf{D}, \mathsf{P}, \mathsf{R}, \mathsf{X}
  \}
\]
where $\mathsf{G}$ = Graph, $\mathsf{V}$ = Vector, $\mathsf{T}$ = Tensor,
$\mathsf{S}$ = Semantic, $\mathsf{D}$ = Document, $\mathsf{P}$ = Temporal,
$\mathsf{R}$ = Provenance, and $\mathsf{X}$ = Spatial.
\end{definition}

Each modality $m \in \modset$ has an associated \emph{modality space} $\Sigma_m$:

\begin{itemize}[leftmargin=*]
  \item $\Sigma_\mathsf{G}$: the set of all labelled directed multigraphs
    (RDF triples or property-graph subgraphs).
  \item $\Sigma_\mathsf{V}$: $\mathbb{R}^d$ for a fixed embedding dimension $d$.
  \item $\Sigma_\mathsf{T}$: the set of typed multi-dimensional arrays
    $\bigcup_{n \geq 1} \mathbb{R}^{d_1 \times \cdots \times d_n}$.
  \item $\Sigma_\mathsf{S}$: the set of CBOR-encoded semantic annotations,
    comprising type URIs and proof blobs.
  \item $\Sigma_\mathsf{D}$: the set of structured text documents with
    inverted-index metadata.
  \item $\Sigma_\mathsf{P}$: the set of Merkle-tree version histories over
    octad snapshots.
  \item $\Sigma_\mathsf{R}$: the set of SHA-256 hash chains recording
    provenance events (creation, transformation, access).
  \item $\Sigma_\mathsf{X}$: the set of geometric objects in a WGS84
    coordinate reference system, indexed by R-tree.
\end{itemize}

\begin{definition}[Octad]\label{def:octad}
An \emph{octad} is a tuple $\octad = (\mathit{id}, \varphi)$ where:
\begin{itemize}[leftmargin=*]
  \item $\mathit{id} \in \mathsf{UUID}$ is a universally unique 128-bit identifier.
  \item $\varphi : \modset \to \bigcup_{m \in \modset} (\Sigma_m \cup \{\bot\})$
    is a \emph{modality function} such that for each $m \in \modset$,
    $\varphi(m) \in \Sigma_m \cup \{\bot\}$, where $\bot$ denotes the absence
    of a representation in modality $m$.
\end{itemize}
\end{definition}

The modality function $\varphi$ maps each of the eight modalities to either a
concrete value in the corresponding modality space or the distinguished value
$\bot$.  An octad may have $\bot$ in some modalities (e.g., an entity without
spatial coordinates), but its identity persists regardless.

\begin{definition}[Octad Store]\label{def:store}
An \emph{octad store} $\mathcal{S}$ is a partial function
$\mathcal{S} : \mathsf{UUID} \rightharpoonup \octad$ together with:
\begin{enumerate}[leftmargin=*]
  \item A \emph{modality index} $\mathcal{I}_m$ for each $m \in \modset$,
    providing efficient retrieval within that modality space.
  \item A \emph{drift monitor} $\driftfn$ that continuously evaluates
    cross-modal coherence (\cref{sec:drift}).
  \item A \emph{repair engine} $\repairfn$ that resolves detected drift
    according to policy (\cref{sec:repair}).
\end{enumerate}
\end{definition}

\subsection{Identity Invariants}

The octad model enforces two fundamental invariants:

\begin{invariant}[Identity Uniqueness]\label{inv:unique}
For any octad store $\mathcal{S}$ and identifiers $\mathit{id}_1 \neq \mathit{id}_2$,
$\mathcal{S}(\mathit{id}_1)$ and $\mathcal{S}(\mathit{id}_2)$ share no modality
representations.  That is, if $\varphi_1(m) \neq \bot$ and $\varphi_2(m) \neq \bot$
for some $m$, then $\varphi_1(m) \neq \varphi_2(m)$.
\end{invariant}

\begin{invariant}[Identity Persistence]\label{inv:persist}
An octad's identifier $\mathit{id}$ is immutable.  Mutations to $\varphi$ do
not change $\mathit{id}$; the temporal modality $\varphi(\mathsf{P})$ records the
full history of all mutations.
\end{invariant}

These invariants guarantee that an entity's identity is stable and
unambiguous across all modalities and across time.

\subsection{Modality Algebra}\label{sec:modality-algebra}

We define algebraic operations over octads to support composition, projection,
and merging.

\begin{definition}[Projection]\label{def:projection}
For an octad $\octad = (\mathit{id}, \varphi)$ and a subset
$M \subseteq \modset$, the \emph{projection} $\pi_M(\octad)$ is the tuple
$(\mathit{id}, \varphi|_M)$ where $\varphi|_M(m) = \varphi(m)$ if $m \in M$,
and $\varphi|_M(m) = \bot$ otherwise.
\end{definition}

\begin{definition}[Merge]\label{def:merge}
Given two octads $\octad_1 = (\mathit{id}, \varphi_1)$ and
$\octad_2 = (\mathit{id}, \varphi_2)$ sharing the same identifier, their
\emph{merge} $\octad_1 \oplus \octad_2 = (\mathit{id}, \varphi_\oplus)$ is
defined by:
\[
  \varphi_\oplus(m) = \begin{cases}
    \varphi_1(m) & \text{if } \varphi_2(m) = \bot \\
    \varphi_2(m) & \text{if } \varphi_1(m) = \bot \\
    \mathsf{resolve}(\varphi_1(m), \varphi_2(m)) & \text{otherwise}
  \end{cases}
\]
where $\mathsf{resolve}$ is a policy-dependent conflict resolution function
(\cref{sec:repair}).
\end{definition}

\begin{proposition}[Merge Commutativity]
If $\mathsf{resolve}$ is commutative, then $\oplus$ is commutative:
$\octad_1 \oplus \octad_2 = \octad_2 \oplus \octad_1$.
\end{proposition}

\begin{proof}
Follows directly from the symmetry of the case analysis in
\cref{def:merge} and the commutativity assumption on $\mathsf{resolve}$.
\end{proof}

\begin{definition}[Enrichment]\label{def:enrichment}
An \emph{enrichment} of octad $\octad = (\mathit{id}, \varphi)$ at modality
$m$ is the operation $\mathsf{enrich}_m(\octad, v)$ producing
$(\mathit{id}, \varphi')$ where $\varphi'(m) = v \in \Sigma_m$ and
$\varphi'(m') = \varphi(m')$ for all $m' \neq m$.  Every enrichment appends
a record to $\varphi'(\mathsf{P})$ and $\varphi'(\mathsf{R})$.
\end{definition}

This enrichment operation is the fundamental write primitive.  It guarantees
that every modality mutation is reflected in both the temporal and provenance
modalities, maintaining the audit trail invariant.

\subsection{Cross-Modal Coherence Constraints}

Beyond identity invariants, the octad model supports user-defined
\emph{coherence constraints} that express expected relationships between
modalities.

\begin{definition}[Coherence Constraint]\label{def:coherence}
A \emph{coherence constraint} is a predicate
$c : \Sigma_{m_1} \times \Sigma_{m_2} \to \{0, 1\}$ for modalities
$m_1, m_2 \in \modset$ that must hold whenever both
$\varphi(m_1) \neq \bot$ and $\varphi(m_2) \neq \bot$.
\end{definition}

Examples of coherence constraints include:
\begin{itemize}[leftmargin=*]
  \item The vector embedding $\varphi(\mathsf{V})$ must be derivable from
    the document content $\varphi(\mathsf{D})$ via a specified embedding
    function $f$: $\|\varphi(\mathsf{V}) - f(\varphi(\mathsf{D}))\|_2 < \epsilon$.
  \item The graph modality $\varphi(\mathsf{G})$ must contain exactly the
    relations declared in the semantic modality $\varphi(\mathsf{S})$.
  \item The spatial modality $\varphi(\mathsf{X})$ must be consistent with
    any geographic references in the document modality $\varphi(\mathsf{D})$.
\end{itemize}

Violation of a coherence constraint constitutes \emph{drift}, which we
formalise in the next section.

% ── 4. Drift Detection and Classification ────────────────────────────────────
\section{Drift Detection and Classification}\label{sec:drift}

\subsection{Drift as a First-Class Concept}

Traditional database systems treat inconsistency as an error condition to be
prevented through transactions, locks, or consensus protocols.  In federated
multi-modal settings, this approach is both impractical and undesirable.
Modalities evolve at different rates: a document may be updated before its
embedding is recomputed; a spatial coordinate may be refined independently of
the underlying graph structure; a provenance chain may lag behind the current
state of the data.

\verisimdb adopts a fundamentally different stance: \emph{drift is not an
error; it is an observable quantity that can be measured, classified, and
managed according to policy.}  This design choice is motivated by real-world
federated environments where strict consistency would require either
(a)~centralised coordination, which violates sovereignty, or (b)~blocking
updates until all modalities are synchronised, which is impractical for
large-scale systems.

\subsection{Formal Drift Metrics}

We define drift as a function of the divergence between modality
representations of the same octad.

\begin{definition}[Pairwise Drift]\label{def:pairwise-drift}
For modalities $m_1, m_2 \in \modset$ and an octad $\octad = (\mathit{id}, \varphi)$,
the \emph{pairwise drift} is:
\[
  \driftfn_{m_1, m_2}(\octad) = d_{m_1, m_2}(\varphi(m_1), \varphi(m_2))
\]
where $d_{m_1, m_2} : \Sigma_{m_1} \times \Sigma_{m_2} \to \mathbb{R}_{\geq 0}$
is a modality-pair-specific distance function, with $d_{m_1, m_2}(\cdot, \bot) =
d_{m_1, m_2}(\bot, \cdot) = 0$ by convention (absent modalities do not drift).
\end{definition}

\begin{definition}[Aggregate Drift]\label{def:agg-drift}
The \emph{aggregate drift} of an octad is the weighted sum over all modality
pairs:
\[
  \driftfn(\octad) = \sum_{\{m_1, m_2\} \subseteq \modset}
    w_{m_1, m_2} \cdot \driftfn_{m_1, m_2}(\octad)
\]
where $w_{m_1, m_2} \geq 0$ are policy-configured weights satisfying
$\sum w_{m_1, m_2} = 1$.
\end{definition}

For eight modalities, there are $\binom{8}{2} = 28$ distinct modality pairs.
In practice, not all pairs have meaningful distance functions; we define
metrics for the most important pairs.

\subsection{Modality-Pair Distance Functions}\label{sec:drift-metrics}

\paragraph{Vector--Document drift ($d_{\mathsf{V},\mathsf{D}}$).}
Given an embedding function $f : \Sigma_\mathsf{D} \to \mathbb{R}^d$:
\[
  d_{\mathsf{V},\mathsf{D}}(\vec{v}, \mathit{doc}) =
    1 - \frac{\vec{v} \cdot f(\mathit{doc})}{\|\vec{v}\|\;\|f(\mathit{doc})\|}
\]
This is the cosine distance between the stored embedding and the embedding
that \emph{would} be computed from the current document content.  A drift
value near 0 indicates that the embedding is fresh; a value near 1 indicates
that the document has changed substantially since the embedding was computed.

\paragraph{Graph--Semantic drift ($d_{\mathsf{G},\mathsf{S}}$).}
Let $E_\mathsf{G}$ be the set of typed edges in the graph modality and
$T_\mathsf{S}$ the set of type assertions in the semantic modality.  Define:
\[
  d_{\mathsf{G},\mathsf{S}}(g, s) =
    1 - \frac{|E_\mathsf{G} \cap T_\mathsf{S}|}{|E_\mathsf{G} \cup T_\mathsf{S}|}
\]
This is the Jaccard distance between the structural assertions in the graph
and the ontological type assertions in the semantic modality.

\paragraph{Temporal--Provenance drift ($d_{\mathsf{P},\mathsf{R}}$).}
Let $H_\mathsf{P}$ be the Merkle tree of version history and $C_\mathsf{R}$
the SHA-256 hash chain of provenance events.  Define:
\[
  d_{\mathsf{P},\mathsf{R}}(h, c) =
    \frac{|\mathsf{leaves}(H_\mathsf{P}) \setminus \mathsf{events}(C_\mathsf{R})|}{
      |\mathsf{leaves}(H_\mathsf{P})|}
\]
This measures the fraction of version-history events not accounted for in the
provenance chain --- an indicator of provenance tracking lag.

\paragraph{Tensor--Vector drift ($d_{\mathsf{T},\mathsf{V}}$).}
When the vector modality is a flattened or projected view of the tensor:
\[
  d_{\mathsf{T},\mathsf{V}}(t, \vec{v}) =
    \frac{\|\mathsf{proj}(t) - \vec{v}\|_2}{\|\mathsf{proj}(t)\|_2}
\]
where $\mathsf{proj} : \Sigma_\mathsf{T} \to \mathbb{R}^d$ is the configured
projection function.

\paragraph{Document--Spatial drift ($d_{\mathsf{D},\mathsf{X}}$).}
Using named-entity recognition to extract geographic references $G_\mathsf{D}$
from the document and comparing with the spatial modality's coordinates
$p_\mathsf{X}$:
\[
  d_{\mathsf{D},\mathsf{X}}(\mathit{doc}, p) =
    \min_{g \in G_\mathsf{D}} \mathsf{haversine}(g, p_\mathsf{X})
\]
normalised to $[0, 1]$ by a configurable maximum distance threshold.

\subsection{Drift Classification}\label{sec:drift-class}

Not all drift is equal.  We classify detected drift into four categories
based on its source and severity:

\begin{definition}[Drift Classes]\label{def:drift-classes}
\begin{enumerate}[leftmargin=*]
  \item \textbf{Representational drift} occurs when a derived modality
    becomes stale relative to its source (e.g., an embedding not yet
    recomputed after a document update).  This is typically benign and
    automatically repairable.
  \item \textbf{Semantic drift} occurs when the meaning of an entity
    changes in one modality but not others (e.g., a retracted paper whose
    graph edges still assert validity).  This requires domain-aware repair.
  \item \textbf{Structural drift} occurs when the topology of
    relationships in the graph modality diverges from the structure implied
    by other modalities (e.g., a document describes a citation that has
    no corresponding graph edge).
  \item \textbf{Provenance drift} occurs when the audit trail in the
    provenance or temporal modalities fails to account for observed changes
    in other modalities.  This may indicate tampering or system failure.
\end{enumerate}
\end{definition}

The classification is computed by a rule engine that evaluates which modality
pairs exhibit drift and maps the pattern to a drift class.  For instance,
drift exclusively in $d_{\mathsf{V},\mathsf{D}}$ is classified as
representational, while drift in $d_{\mathsf{G},\mathsf{S}}$ with concurrent
drift in $d_{\mathsf{P},\mathsf{R}}$ is classified as provenance drift
(indicating an untracked semantic change).

\subsection{Drift Detection Algorithm}

\begin{algorithm}[t]
\caption{Periodic Drift Detection}\label{alg:drift-detect}
\begin{algorithmic}[1]
\Require Octad store $\mathcal{S}$, sampling interval $\tau$,
  thresholds $\theta_\mathsf{soft}, \theta_\mathsf{hard}$
\Loop
  \State $\mathit{sample} \gets \mathsf{random\_sample}(\mathcal{S}, k)$
  \For{$\octad \in \mathit{sample}$}
    \For{$(m_1, m_2) \in \mathsf{active\_pairs}(\octad)$}
      \State $\delta \gets \driftfn_{m_1,m_2}(\octad)$
      \If{$\delta > \theta_\mathsf{hard}$}
        \State $\mathsf{emit\_alert}(\octad, m_1, m_2, \mathsf{HARD})$
      \ElsIf{$\delta > \theta_\mathsf{soft}$}
        \State $\mathsf{mark\_drifted}(\octad, m_1, m_2)$
      \EndIf
    \EndFor
  \EndFor
  \State $\mathsf{sleep}(\tau)$
\EndLoop
\end{algorithmic}
\end{algorithm}

\Cref{alg:drift-detect} shows the core drift detection loop.  The algorithm
operates by random sampling to amortise the cost of cross-modal comparisons
over large stores.  The function $\mathsf{active\_pairs}(\octad)$ returns
only those modality pairs for which both modalities are non-$\bot$ and a
distance function is defined.  Two thresholds govern the response: a soft
threshold triggers marking (the octad is flagged but no immediate action is
taken), and a hard threshold triggers an alert that initiates repair.

The sampling rate $k$ and interval $\tau$ are tunable: aggressive settings
($k = |\mathcal{S}|$, $\tau = 0$) provide continuous full-scan monitoring
at the cost of throughput, while conservative settings enable drift detection
with minimal overhead.

% ── 5. Policy-Driven Repair ──────────────────────────────────────────────────
\section{Policy-Driven Repair}\label{sec:repair}

\subsection{Repair Strategies}

Once drift is detected and classified, the system must determine how to
resolve it.  We define three \emph{automation levels}:

\begin{definition}[Repair Automation Levels]
\begin{enumerate}[leftmargin=*]
  \item \textbf{Automatic repair.}  The system identifies the
    \emph{authoritative modality} for the drifted pair and regenerates
    the non-authoritative modality.  Example: re-embedding a document
    when $d_{\mathsf{V},\mathsf{D}}$ exceeds threshold.
  \item \textbf{Semi-automatic repair.}  The system proposes a repair
    action and awaits confirmation from a designated \emph{domain custodian}.
    Example: updating graph edges when a semantic annotation changes.
  \item \textbf{Manual repair.}  The system logs the drift event and
    creates a review ticket.  Example: resolving conflicting provenance
    chains that may indicate data tampering.
\end{enumerate}
\end{definition}

\subsection{The Repair Policy Language}

Repair policies are expressed as rules mapping drift class $\times$ severity
to an automation level and a repair action.  Formally:

\begin{definition}[Repair Policy]\label{def:repair-policy}
A \emph{repair policy} is a function:
\[
  \policyfn : \mathsf{DriftClass} \times \mathbb{R}_{\geq 0} \to
    \mathsf{AutoLevel} \times \mathsf{Action}
\]
where $\mathsf{DriftClass}$ is one of the four classes in
\cref{def:drift-classes}, $\mathsf{AutoLevel} \in \{
\mathsf{auto}, \mathsf{semi}, \mathsf{manual}\}$, and $\mathsf{Action}$
specifies the concrete repair operation.
\end{definition}

Policies are encoded in CBOR and stored in the semantic modality of a
designated \emph{policy octad}, itself subject to the same identity and
versioning invariants as any other octad.  This design ensures that policy
changes are auditable and versioned.

\subsection{Authoritative Modality Selection}

When automatic repair is triggered, the system must determine which modality
is \emph{authoritative} --- i.e., which modality's current state should be
treated as ground truth.

\begin{definition}[Authority Function]\label{def:authority}
An \emph{authority function} $\alpha : \modset \times \modset \to \modset$
maps each modality pair to the authoritative member.  The default authority
ordering is:
\[
  \mathsf{D} > \mathsf{G} > \mathsf{S} > \mathsf{V} > \mathsf{T} >
  \mathsf{X} > \mathsf{P} > \mathsf{R}
\]
where $m_1 > m_2$ means $m_1$ is authoritative over $m_2$.
\end{definition}

The rationale for this ordering is that human-authored modalities (documents,
graphs, semantic annotations) take precedence over derived modalities
(embeddings, tensors) and infrastructure modalities (temporal, provenance).
The ordering is configurable per octad store to accommodate domain-specific
requirements.

\subsection{Repair Execution}

\begin{algorithm}[t]
\caption{Policy-Driven Repair Execution}\label{alg:repair}
\begin{algorithmic}[1]
\Require Drifted octad $\octad$, modality pair $(m_1, m_2)$,
  drift class $c$, drift value $\delta$
\State $(\mathit{level}, \mathit{action}) \gets \policyfn(c, \delta)$
\If{$\mathit{level} = \mathsf{auto}$}
  \State $m_\alpha \gets \alpha(m_1, m_2)$
  \State $m_\beta \gets (m_1, m_2) \setminus \{m_\alpha\}$
  \State $v' \gets \mathsf{derive}(\varphi(m_\alpha), m_\beta)$
  \State $\octad' \gets \mathsf{enrich}_{m_\beta}(\octad, v')$
  \State $\mathsf{commit}(\octad')$
\ElsIf{$\mathit{level} = \mathsf{semi}$}
  \State $\mathsf{propose\_repair}(\octad, m_1, m_2, \mathit{action})$
  \State \textbf{await} custodian approval or timeout
\Else
  \State $\mathsf{log\_for\_review}(\octad, m_1, m_2, c, \delta)$
\EndIf
\end{algorithmic}
\end{algorithm}

The $\mathsf{derive}$ function in \cref{alg:repair} is modality-pair-specific:
for $(\mathsf{D}, \mathsf{V})$, it invokes the embedding model; for
$(\mathsf{G}, \mathsf{S})$, it extracts typed edges from semantic annotations.
Every repair action, regardless of automation level, appends an entry to both
the temporal and provenance modalities.

\subsection{Immutable Audit Trail}

All repair actions produce an audit record:
\[
  \mathsf{AuditRecord} = (
    \mathit{timestamp},\;
    \mathit{octad\_id},\;
    \mathit{drift\_class},\;
    \mathit{drift\_value},\;
    \mathit{repair\_action},\;
    \mathit{actor},\;
    \mathit{prev\_hash}
  )
\]

Records form a hash chain: $\mathit{hash}_i = \mathsf{SHA256}(
\mathit{AuditRecord}_i \| \mathit{hash}_{i-1})$.  This chain is stored in
the provenance modality and is verifiable by any federated participant.

% ── 6. Cross-Modal Query ─────────────────────────────────────────────────────
\section{Cross-Modal Query}\label{sec:query}

\subsection{Query Model}

A key advantage of the octad model is that queries can \emph{span modality
boundaries} while maintaining type safety.  We define a cross-modal query
language, VQL (VeriSim Query Language), that supports modality-specific
operations composed through a common identity layer.

\begin{definition}[Cross-Modal Query]\label{def:query}
A \emph{cross-modal query} $q$ is a composition of modality-specific
query fragments:
\[
  q = \queryfn_{m_1}(\mathit{pred}_1) \bowtie_{\mathit{id}}
      \queryfn_{m_2}(\mathit{pred}_2) \bowtie_{\mathit{id}} \cdots
      \bowtie_{\mathit{id}} \queryfn_{m_k}(\mathit{pred}_k)
\]
where $\queryfn_{m_i}(\mathit{pred}_i)$ returns a set of octad identifiers
matching predicate $\mathit{pred}_i$ in modality $m_i$, and
$\bowtie_{\mathit{id}}$ is the natural join on octad identity.
\end{definition}

This formulation enables queries such as ``find all entities that are within
10km of London (spatial), were authored after 2020 (temporal), cite a
retracted paper (graph), and have a cosine similarity $> 0.9$ to a query
embedding (vector).''

\subsection{Query Planning}

The cross-modal query planner operates in three phases:

\begin{enumerate}[leftmargin=*]
  \item \textbf{Decomposition.}  The query is parsed into modality-specific
    fragments.  Each fragment is routed to the appropriate modality index.
  \item \textbf{Selectivity estimation.}  The planner estimates the
    selectivity of each fragment using modality-specific statistics
    (cardinality estimates for graph patterns, distance distributions
    for vector queries, spatial density for R-tree bounds).
  \item \textbf{Execution ordering.}  Fragments are executed in order of
    decreasing selectivity (most selective first) to minimise the
    intermediate result set.  The identity join is implemented as a
    hash-based intersection of UUID sets.
\end{enumerate}

\subsection{VQL Syntax (Abbreviated)}

\begin{lstlisting}[language=SQL,caption={Cross-modal VQL query example}]
SELECT octad.id, octad.document.title,
       octad.spatial.coordinates
FROM octads
WHERE octad.vector NEAR [0.1, 0.3, ...] LIMIT 100
  AND octad.graph MATCHES
      (?x :cites ?y WHERE ?y :status "retracted")
  AND octad.spatial WITHIN radius(51.5, -0.12, 10km)
  AND octad.temporal.created > "2020-01-01"
ORDER BY octad.vector.similarity DESC
\end{lstlisting}

The query in Listing~1 demonstrates the four-modality span: vector similarity
(line 4), graph pattern matching (lines 5--6), spatial containment (line 7),
and temporal filtering (line 8).  The result includes document and spatial
projections of matching octads.

\subsection{Type Safety}

VQL is statically typed: each modality accessor (e.g., \texttt{octad.vector},
\texttt{octad.graph}) returns a value of the corresponding modality type.
The type checker rejects queries that apply modality-inappropriate operations
(e.g., applying \texttt{NEAR} to a document modality or \texttt{MATCHES} to
a vector modality).  This type discipline is enforced at parse time, before
any query execution.

% ── 7. Federated Verification ────────────────────────────────────────────────
\section{Federated Verification}\label{sec:federation}

\subsection{Sovereignty Model}

In a federated deployment, each participating institution operates its own
octad store.  \verisimdb does not assume a central coordinator or a shared
consensus protocol.  Instead, federation is achieved through:

\begin{enumerate}[leftmargin=*]
  \item \textbf{Shared identity namespace.}  All federated stores use the
    same UUID generation scheme, ensuring that an octad created at
    institution A can be referenced by institution B without identifier
    collision.
  \item \textbf{Bilateral verification.}  Two stores that wish to
    interoperate exchange \emph{drift attestations}: signed statements of
    the form ``octad $\mathit{id}$ has aggregate drift $\driftfn(\octad)
    \leq \theta$ as of timestamp $t$, attested by store $S$.''
  \item \textbf{Trust policies.}  Each store independently configures
    which remote attestations it accepts, from which attestors, and
    subject to what freshness requirements.
\end{enumerate}

\subsection{Federation Protocol}

\begin{definition}[Drift Attestation]\label{def:attestation}
A \emph{drift attestation} is a signed tuple:
\[
  A = (
    \mathit{octad\_id},\;
    \driftfn(\octad),\;
    t,\;
    \mathit{store\_id},\;
    \mathsf{sign}(\mathit{octad\_id} \| \driftfn(\octad) \| t,\; \mathit{sk})
  )
\]
where $\mathit{sk}$ is the attesting store's private key and $t$ is the
attestation timestamp.
\end{definition}

The federation protocol proceeds as follows:

\begin{enumerate}[leftmargin=*]
  \item \textbf{Discovery.}  Stores publish their public keys and endpoint
    URIs to a shared registry (or discover them via DNS-based service
    discovery).
  \item \textbf{Attestation exchange.}  When a store receives a query
    referencing a remote octad, it requests a drift attestation from the
    remote store.
  \item \textbf{Verification.}  The requesting store verifies the
    attestation signature and checks that $\driftfn(\octad) \leq \theta$
    where $\theta$ is the local trust policy's maximum acceptable drift.
  \item \textbf{Conditional access.}  If verification succeeds, the
    requesting store proxies the query to the remote store.  If it fails,
    the query returns a \emph{drift warning} to the client, who may
    choose to proceed with degraded trust or abort.
\end{enumerate}

This protocol is deliberately lightweight: it adds one round-trip (the
attestation request) to cross-store queries, and the attestation itself is
a compact signed message (approximately 200 bytes).

\subsection{Conflict Resolution in Federation}

When the same octad is hosted by multiple federated stores (replicated
octads), conflicts may arise.  \verisimdb uses a \emph{last-writer-wins}
strategy by default, with the temporal modality providing a total order on
writes.  Stores may optionally employ application-specific merge functions
via the $\mathsf{resolve}$ function in \cref{def:merge}.

\begin{theorem}[Attestation Soundness]\label{thm:attestation}
If a drift attestation $A$ is valid (signature verification succeeds and
$t$ is within the freshness window), then the attesting store's local
drift measurement of the referenced octad was at most $\driftfn(\octad)$
at time $t$.
\end{theorem}

\begin{proof}
Follows from the unforgeability of the signature scheme and the
deterministic computation of $\driftfn$.  The attesting store computes
$\driftfn(\octad)$ at time $t$, signs the result, and publishes it.
A verifier who accepts the signature knows that the attesting store
produced this exact drift value.  The freshness window guards against
replay of stale attestations.
\end{proof}

% ── 8. Implementation ────────────────────────────────────────────────────────
\section{Implementation}\label{sec:implementation}

\subsection{Architecture Overview}

\verisimdb is implemented as a two-tier system: a \emph{Rust core} providing
modality-specific storage and query engines, and an \emph{Elixir orchestration
layer} providing distributed coordination, drift monitoring, and query
routing.  The two tiers communicate via HTTP and a C-ABI foreign function
interface for performance-critical paths.

\begin{table}[t]
\centering
\caption{Modality implementation summary}\label{tab:impl}
\begin{tabularx}{\columnwidth}{@{}lXl@{}}
\toprule
\textbf{Modality} & \textbf{Backend} & \textbf{Crate} \\
\midrule
Graph     & Oxigraph (RDF/SPARQL)      & \texttt{verisim-graph} \\
Vector    & Custom HNSW                & \texttt{verisim-vector} \\
Tensor    & \texttt{ndarray} / Burn    & \texttt{verisim-tensor} \\
Semantic  & CBOR (via \texttt{ciborium}) & \texttt{verisim-semantic} \\
Document  & Tantivy (inverted index)   & \texttt{verisim-document} \\
Temporal  & Merkle-tree snapshots      & \texttt{verisim-temporal} \\
Provenance & SHA-256 hash chains       & \texttt{verisim-provenance} \\
Spatial   & R-tree (\texttt{rstar})    & \texttt{verisim-spatial} \\
\bottomrule
\end{tabularx}
\end{table}

\Cref{tab:impl} summarises the backend technology for each modality.
All crates are written in safe Rust, with \texttt{unsafe} blocks limited
to FFI boundaries.

\subsection{Rust Core}

The Rust core is organised as a Cargo workspace with one crate per modality,
plus cross-cutting crates for the octad abstraction (\texttt{verisim-octad}),
drift detection (\texttt{verisim-drift}), and self-normalisation
(\texttt{verisim-normalizer}).

\paragraph{Graph modality.}
Oxigraph provides an embedded RDF store with SPARQL 1.1 support.
\verisimdb extends Oxigraph with property-graph semantics by encoding
property-graph edges as reified RDF statements, enabling a single storage
backend to support both RDF and property-graph query patterns.  The graph
modality supports both in-memory and persistent (file-backed) modes.

\paragraph{Vector modality.}
The HNSW index is implemented with configurable parameters
($M = 16$, $\mathit{ef}_\mathsf{construction} = 200$ by default).
Vectors are stored in memory-mapped files, enabling datasets larger than
available RAM.  Distance functions include cosine similarity, Euclidean
distance, and inner product.

\paragraph{Tensor modality.}
Multi-dimensional arrays are stored using the \texttt{ndarray} crate for
CPU operations, with optional Burn integration for GPU-accelerated tensor
computations.  Tensors are serialised in a custom binary format with
dimension metadata.

\paragraph{Semantic modality.}
Semantic annotations are CBOR-encoded structures containing type URIs
(linking to external ontologies), proof blobs (for formally verified
properties), and policy references.  The \texttt{ciborium} crate provides
CBOR serialisation with schema validation.

\paragraph{Document modality.}
Tantivy provides full-text indexing with BM25 ranking, phrase queries,
and boolean operators.  Documents are stored with field-level schema
(title, body, metadata) and support incremental index updates.

\paragraph{Temporal modality.}
Version history is maintained as a Merkle tree of octad snapshots.
Each node in the tree contains a cryptographic hash of the octad state
at that version, enabling efficient verification of version integrity
and time-travel queries of the form ``retrieve the octad as it existed
at timestamp $t$.''

\paragraph{Provenance modality.}
The provenance chain is a SHA-256 linked list of events.  Each event
records the actor, action, timestamp, and the hash of the previous event.
The chain supports efficient verification: given any event, the entire
history back to the genesis event can be verified in $O(n)$ hash
computations.

\paragraph{Spatial modality.}
The R-tree index (via the \texttt{rstar} crate) supports WGS84
coordinates with point, bounding-box, and polygon geometries.
Supported queries include radius search, bounding-box containment,
and $k$-nearest-neighbour.

\subsection{Elixir Orchestration Layer}

The Elixir layer is built on OTP (Open Telecom Platform) and provides:

\begin{itemize}[leftmargin=*]
  \item \textbf{Entity servers.}  Each octad is managed by a GenServer
    process that coordinates access across modalities.  OTP supervision
    ensures fault tolerance: if an entity server crashes, it is
    automatically restarted with its last known state.
  \item \textbf{Drift monitor.}  A GenStage pipeline continuously samples
    octads and evaluates drift metrics.  The pipeline consists of three
    stages: \emph{sampler} (selects octads), \emph{evaluator} (computes
    drift), and \emph{dispatcher} (routes alerts).
  \item \textbf{Query router.}  Incoming queries are decomposed into
    modality-specific fragments, dispatched to the appropriate Rust
    crate via HTTP or NIF, and the results are joined on octad identity.
  \item \textbf{Schema registry.}  Maintains the mapping between octad
    types and their expected modality configurations, coherence
    constraints, and repair policies.
\end{itemize}

\subsection{Persistence and Recovery}

\verisimdb supports two deployment modes:

\begin{enumerate}[leftmargin=*]
  \item \textbf{In-memory mode.}  All modality stores reside in RAM.
    Suitable for testing, ephemeral workloads, and embedded deployments.
  \item \textbf{Persistent mode.}  Modality stores are backed by files
    (memory-mapped where possible).  A write-ahead log (WAL) ensures
    durability: all enrichment operations are logged before being applied
    to the modality stores.  Recovery replays the WAL from the last
    checkpoint.
\end{enumerate}

\subsection{Container Deployment}

\verisimdb is packaged as an OCI-compliant container image built on
Chainguard base images for supply-chain security.  The container includes
the Rust core, Elixir orchestration layer, and a reverse proxy for TLS
termination and rate limiting.

% ── 9. Evaluation ────────────────────────────────────────────────────────────
\section{Evaluation}\label{sec:evaluation}

We evaluate \verisimdb along three axes: drift detection accuracy,
cross-modal query performance, and federation overhead.

\subsection{Experimental Setup}

All experiments were conducted on a single server with an AMD Ryzen 9 7950X
(16 cores / 32 threads), 64 GB DDR5 RAM, and an NVMe SSD (Samsung 990 Pro).
The Rust core was compiled with \texttt{--release} optimisations (LTO enabled).
The Elixir layer ran on BEAM/OTP 27.

\paragraph{Datasets.}
We use two datasets:
\begin{itemize}[leftmargin=*]
  \item \textbf{SynthOctad-1M.}  A synthetic dataset of 1 million octads
    with all eight modalities populated.  Documents are generated from
    Wikipedia abstracts; embeddings from a 384-dimensional MiniLM model;
    graph edges from Wikidata; spatial coordinates uniformly distributed
    over the globe; temporal histories with 5--20 versions per octad.
  \item \textbf{GeoScholar-100K.}  A real-world dataset of 100{,}000
    academic publications with geographic metadata, citation graphs from
    OpenAlex, and embeddings from Specter2.  Provenance chains are
    derived from CrossRef event data.
\end{itemize}

\subsection{Drift Detection Accuracy}

To evaluate drift detection, we introduce controlled drift into the
SynthOctad-1M dataset by independently modifying modalities without updating
their counterparts.  We inject three drift types:

\begin{enumerate}[leftmargin=*]
  \item \textbf{Embedding staleness.}  We update 10\% of documents without
    recomputing embeddings.
  \item \textbf{Graph--semantic divergence.}  We add graph edges to 5\% of
    octads without updating semantic annotations.
  \item \textbf{Provenance gaps.}  We modify 3\% of octads through a
    side channel that bypasses provenance logging.
\end{enumerate}

\begin{table}[t]
\centering
\caption{Drift detection accuracy on SynthOctad-1M}\label{tab:drift-accuracy}
\begin{tabular}{@{}lccc@{}}
\toprule
\textbf{Drift Type} & \textbf{Precision} & \textbf{Recall} & \textbf{F1} \\
\midrule
Embedding staleness      & 0.97 & 0.94 & 0.955 \\
Graph--semantic div.     & 0.93 & 0.89 & 0.910 \\
Provenance gaps          & 1.00 & 0.98 & 0.990 \\
\midrule
\textbf{Macro average}   & 0.97 & 0.94 & 0.952 \\
\bottomrule
\end{tabular}
\end{table}

\Cref{tab:drift-accuracy} shows the results.  Embedding staleness is detected
with high precision because cosine distance provides a reliable signal when
documents change.  Graph--semantic divergence is harder to detect precisely
because small structural changes (e.g., adding a single edge) produce small
Jaccard distance changes that may fall below the soft threshold.  Provenance
gaps are detected with near-perfect accuracy because the hash-chain structure
makes any missing event immediately evident.

\subsection{Cross-Modal Query Performance}

We measure end-to-end latency for cross-modal queries spanning 1, 2, 4,
and 8 modalities on the SynthOctad-1M dataset.

\begin{table}[t]
\centering
\caption{Cross-modal query latency (ms), median of 1000 runs}\label{tab:query-latency}
\begin{tabular}{@{}lcccc@{}}
\toprule
\textbf{Modalities} & \textbf{p50} & \textbf{p95} & \textbf{p99} & \textbf{Throughput} \\
 & (ms) & (ms) & (ms) & (q/s) \\
\midrule
1 (vector only)     &  1.2 &  3.8 &  7.1 & 12{,}400 \\
2 (vector + graph)  &  4.7 & 11.2 & 19.3 &  5{,}800 \\
4 (V+G+D+X)        & 12.1 & 28.4 & 41.7 &  2{,}100 \\
8 (all modalities)  & 31.8 & 67.2 & 94.5 &    780 \\
\bottomrule
\end{tabular}
\end{table}

\Cref{tab:query-latency} shows that single-modality queries achieve sub-2ms
median latency, while full eight-modality queries complete in under 32ms at
the median.  The latency increase is sub-linear in the number of modalities
because the query planner executes the most selective fragment first and
subsequent fragments operate on progressively smaller candidate sets.

\subsection{Federation Overhead}

We measure the overhead of federated verification by comparing query latency
for local vs.\ federated queries on GeoScholar-100K, simulating two stores
on the same network with 0.5ms RTT.

\begin{table}[t]
\centering
\caption{Federation overhead (ms, median)}\label{tab:federation}
\begin{tabular}{@{}lcc@{}}
\toprule
\textbf{Query Type} & \textbf{Local} & \textbf{Federated} \\
\midrule
Single-modality   &  1.8 &   3.1 \\
Two-modality      &  5.2 &   7.8 \\
Four-modality     & 13.4 &  18.2 \\
\bottomrule
\end{tabular}
\end{table}

\Cref{tab:federation} shows that federation adds approximately 1.3--4.8ms
per query, dominated by the attestation round-trip and signature verification.
The Ed25519 signature verification takes approximately 0.06ms; the remainder
is network latency.

\subsection{Scalability}

We measure index build time and memory consumption as a function of octad
count on SynthOctad-1M.

\begin{table}[t]
\centering
\caption{Scalability metrics}\label{tab:scalability}
\begin{tabular}{@{}lrrr@{}}
\toprule
\textbf{Octad Count} & \textbf{Build Time} & \textbf{Memory} & \textbf{Disk} \\
\midrule
10K     &   2.1s &  0.4 GB &  0.3 GB \\
100K    &  18.7s &  3.2 GB &  2.8 GB \\
1M      & 194.3s & 28.1 GB & 24.6 GB \\
\bottomrule
\end{tabular}
\end{table}

Build time and memory scale approximately linearly.  The per-octad memory
cost of approximately 28 KB is dominated by the HNSW index (384-dimensional
vectors at 4 bytes per dimension, plus graph overhead for the HNSW layers)
and the Tantivy inverted index.

% ── 10. Related Work ─────────────────────────────────────────────────────────
\section{Related Work}\label{sec:related}

\subsection{Multi-Model Databases}

\textbf{ArangoDB}~\cite{arangodb2023} is a mature multi-model database
supporting document, graph, and key-value access patterns over a single
storage engine (RocksDB).  It provides AQL (ArangoDB Query Language) for
cross-model queries.  However, ArangoDB does not maintain formal coherence
guarantees between models: a document and its corresponding graph
representation can diverge without detection.  ArangoDB also lacks native
vector search and spatial indexing.

\textbf{SurrealDB}~\cite{surrealdb2024} extends the multi-model paradigm
with record links, full-text search, vector similarity, and a flexible
schema system.  SurrealQL provides a unified query interface.  SurrealDB
represents the closest existing system to \verisimdb in ambition, but differs
in three key respects: (1)~it does not define formal drift semantics between
models, (2)~it does not support temporal versioning with Merkle-tree
integrity, and (3)~its federation story relies on traditional
primary-replica patterns rather than sovereign attestation-based verification.

\textbf{TerminusDB}~\cite{terminusdb2023} applies version-control semantics
to graph data, supporting branch, merge, and diff operations.  Its focus on
temporal integrity over graph data makes it complementary to the octad model's
temporal modality, but TerminusDB does not extend versioning to other
modalities (vectors, tensors, documents, spatial data).

\subsection{Decentralised and Federated Data Systems}

\textbf{IPFS}~\cite{benet2014ipfs} provides content-addressable, distributed
storage using Merkle DAGs.  Data is identified by its hash, ensuring
immutability.  However, IPFS is modality-agnostic: it stores opaque blobs
with no cross-modal coherence semantics.  The octad model's provenance
modality is conceptually related to IPFS's content addressing, but extends
it with actor tracking and temporal ordering.

\textbf{Solid}~\cite{sambra2016solid} enables user-controlled data pods with
fine-grained access control using Web Access Control (WAC) and linked data.
Solid's sovereignty model is aligned with \verisimdb's federation goals, but
Solid does not provide cross-modal query capabilities or drift detection.

\textbf{Ceramic}~\cite{ceramic2023} is a decentralised data network built on
IPFS that supports mutable, version-controlled data streams.  It provides
eventual consistency through CRDT-based conflict resolution.  Unlike
\verisimdb, Ceramic treats consistency as a convergence property rather than
a drift-aware policy decision.

\subsection{Knowledge Graph Platforms}

\textbf{Wikidata}~\cite{vrandecic2014wikidata} maintains a collaboratively
edited knowledge graph with provenance tracking (references and qualifiers).
While Wikidata tracks the source of claims, it does not maintain vector,
tensor, or spatial representations of entities, and its consistency model
relies on human editorial review rather than automated drift detection.

\textbf{Neo4j}~\cite{neo4j2023} is the most widely deployed property-graph
database.  Recent versions add vector search capabilities, but this is
implemented as a separate index layer rather than a co-equal modality with
coherence guarantees.

\subsection{Vector Database Systems}

\textbf{Milvus}~\cite{wang2021milvus}, \textbf{Qdrant}, and \textbf{Weaviate}
are purpose-built vector databases that excel at approximate nearest-neighbour
search.  Weaviate adds a module system for integrating with embedding models
and supports hybrid search (vector + keyword).  None of these systems maintain
formal identity across non-vector representations or detect when stored
embeddings diverge from their source data.

\subsection{Drift and Data Quality}

The concept of \emph{data drift} is well-studied in machine learning
operations (MLOps), where it refers to changes in the statistical
distribution of model inputs or outputs~\cite{lu2019learning}.  Systems
such as Evidently AI, Arize, and WhyLabs provide drift monitoring for
ML pipelines.  \verisimdb generalises this concept from the ML-specific
setting to arbitrary multi-modal data, and adds formal metrics, policy-driven
repair, and federation-aware attestation.

% ── 11. Discussion ───────────────────────────────────────────────────────────
\section{Discussion}\label{sec:discussion}

\subsection{Why Eight Modalities?}

The choice of exactly eight modalities is a pragmatic design decision, not
a theoretical necessity.  The octad model is extensible: adding a ninth
modality (e.g., audio waveforms or time-series sensor data) requires
defining its modality space $\Sigma_m$, implementing a storage backend,
and defining distance functions for relevant modality pairs.  The formal
framework in \cref{sec:octad} generalises straightforwardly to $|\modset| = n$.
We chose eight because they cover the representational needs of our target
domains while remaining tractable for drift analysis ($\binom{8}{2} = 28$
pairs).

\subsection{Limitations}

\paragraph{Drift metric design.}  The quality of drift detection depends on
the quality of the distance functions $d_{m_1, m_2}$.  For some modality
pairs (e.g., tensor--spatial), we have not yet identified meaningful distance
functions.  Users must supply domain-specific metrics for uncommon pairs.

\paragraph{Embedding model dependency.}  The vector--document drift metric
assumes a deterministic embedding function.  In practice, embedding models
are updated, and a new model version will produce different embeddings for
the same document, creating apparent drift that is actually model evolution.
We mitigate this by versioning the embedding function identifier alongside
the vector modality.

\paragraph{Scalability ceiling.}  The current implementation stores the HNSW
index in memory-mapped files, which limits vector search to datasets that fit
in virtual address space.  For datasets exceeding several billion vectors,
a distributed HNSW implementation would be required.

\paragraph{Formal verification.}  While we store proof blobs in the semantic
modality, the current system does not verify these proofs at query time.
Integration with a proof checker (e.g., for Idris2 or Lean-generated proofs)
is planned but not yet implemented.

\subsection{Ethical Considerations}

The drift-aware model raises ethical questions about \emph{who decides}
repair policies.  When a contested historical narrative is represented as
an octad, the authority function determines which modality's version of
events prevails.  We address this by making repair policies explicit,
versioned, and auditable: the provenance chain records every policy
decision, and the temporal modality preserves every prior state.  This
does not resolve the underlying ethical question, but it ensures that
the decision-making process is transparent and reversible.

% ── 12. Conclusion ───────────────────────────────────────────────────────────
\section{Conclusion}\label{sec:conclusion}

We have presented the Octad Data Model, a formal framework for multi-modal
data identity in which every entity exists simultaneously across eight
modalities --- Graph, Vector, Tensor, Semantic, Document, Temporal,
Provenance, and Spatial --- bound to a single UUID-based identity.  The
key departure from existing multi-model databases is the treatment of
cross-modal drift as a first-class concept: drift is detected through
formal metrics, classified by type, and resolved through explicit,
auditable policies.

The federated verification protocol enables sovereign data stores to
interoperate without centralised control, using signed drift attestations
to establish cross-institutional trust.  Our experimental evaluation
demonstrates that drift detection achieves F1 scores above 0.91 across all
drift types, cross-modal queries spanning all eight modalities complete in
under 32ms at the median, and federation adds modest overhead (under 5ms)
to query latency.

The octad model addresses a growing need in data-intensive systems: as the
number of representational modalities per entity increases, so does the
risk of silent divergence between representations.  By making this
divergence observable, classifiable, and policy-governed, \verisimdb
provides a principled foundation for multi-modal data management in
federated environments.

\paragraph{Future work.}
Planned extensions include: (1)~integration of formal proof verification
for semantic modality blobs at query time, (2)~distributed HNSW for
billion-scale vector datasets, (3)~a visual drift dashboard for
non-technical domain custodians, and (4)~benchmarking on additional
real-world federated deployments.

\subsection*{Acknowledgements}

The author thanks the open-source communities behind Oxigraph, Tantivy,
and the Rust and Elixir ecosystems for the foundational infrastructure on
which \verisimdb is built.

% ── References ───────────────────────────────────────────────────────────────
\begin{thebibliography}{99}

\bibitem{arangodb2023}
ArangoDB GmbH.
\newblock ArangoDB: A native multi-model database.
\newblock \url{https://www.arangodb.com/}, 2023.

\bibitem{benet2014ipfs}
J.~Benet.
\newblock {IPFS} -- Content addressed, versioned, {P2P} file system.
\newblock \emph{arXiv preprint arXiv:1407.3561}, 2014.

\bibitem{ceramic2023}
Ceramic Network.
\newblock Ceramic: A decentralized data network.
\newblock \url{https://ceramic.network/}, 2023.

\bibitem{jegou2011pq}
H.~J\'{e}gou, M.~Douze, and C.~Schmid.
\newblock Product quantization for nearest neighbor search.
\newblock \emph{IEEE Transactions on Pattern Analysis and Machine Intelligence},
  33(1):117--128, 2011.

\bibitem{lu2019learning}
J.~Lu, A.~Liu, F.~Dong, F.~Gu, J.~Gama, and G.~Zhang.
\newblock Learning under concept drift: A review.
\newblock \emph{IEEE Transactions on Knowledge and Data Engineering},
  31(12):2346--2363, 2019.

\bibitem{lu2019multimodel}
J.~Lu and I.~Holubov\'{a}.
\newblock Multi-model databases: A new journey to handle the variety of data.
\newblock \emph{ACM Computing Surveys}, 52(3):1--38, 2019.

\bibitem{malkov2018hnsw}
Y.~A.~Malkov and D.~A.~Yashunin.
\newblock Efficient and robust approximate nearest neighbor using hierarchical
  navigable small world graphs.
\newblock \emph{IEEE Transactions on Pattern Analysis and Machine Intelligence},
  42(4):824--836, 2020.

\bibitem{neo4j2023}
Neo4j, Inc.
\newblock Neo4j graph database platform.
\newblock \url{https://neo4j.com/}, 2023.

\bibitem{oxigraph2023}
T.~Tanon.
\newblock Oxigraph: An embedded {SPARQL} store.
\newblock \url{https://oxigraph.org/}, 2023.

\bibitem{sambra2016solid}
A.~V.~Sambra, E.~Mansour, S.~Hawke, M.~Zerber, N.~Greco, A.~Ghanem,
  D.~Zagidulin, A.~Aboulnaga, and T.~Berners-Lee.
\newblock Solid: A platform for decentralized social applications based on
  linked data.
\newblock \emph{MIT CSAIL \& Qatar Computing Research Institute, Tech.\ Rep.},
  2016.

\bibitem{snodgrass1999temporal}
R.~T.~Snodgrass.
\newblock \emph{Developing Time-Oriented Database Applications in {SQL}}.
\newblock Morgan Kaufmann, 1999.

\bibitem{surrealdb2024}
SurrealDB Ltd.
\newblock SurrealDB: The ultimate multi-model database.
\newblock \url{https://surrealdb.com/}, 2024.

\bibitem{terminusdb2023}
TerminusDB.
\newblock TerminusDB: Version control for data.
\newblock \url{https://terminusdb.com/}, 2023.

\bibitem{uber2018h3}
Uber Technologies, Inc.
\newblock {H3}: A hexagonal hierarchical geospatial indexing system.
\newblock \url{https://h3geo.org/}, 2018.

\bibitem{vrandecic2014wikidata}
D.~Vrande\v{c}i\'{c} and M.~Kr\"{o}tzsch.
\newblock Wikidata: A free collaborative knowledgebase.
\newblock \emph{Communications of the ACM}, 57(10):78--85, 2014.

\bibitem{w3cprov2013}
{W3C}.
\newblock {PROV-DM}: The {PROV} data model.
\newblock W3C Recommendation, 2013.
\newblock \url{https://www.w3.org/TR/prov-dm/}.

\bibitem{wang2021milvus}
J.~Wang, X.~Yi, R.~Guo, H.~Jin, P.~Xu, S.~Li, X.~Wang, X.~Guo, C.~Li,
  X.~Xu, et~al.
\newblock Milvus: A purpose-built vector data management system.
\newblock In \emph{Proc.\ ACM SIGMOD}, pp.\ 2614--2627, 2021.

\bibitem{wilkinson2016fair}
M.~D.~Wilkinson et~al.
\newblock The {FAIR} guiding principles for scientific data management and
  stewardship.
\newblock \emph{Scientific Data}, 3:160018, 2016.

\end{thebibliography}

% ── Appendix ─────────────────────────────────────────────────────────────────
\appendix

\section{Modality Space Formal Specifications}\label{app:modality-spaces}

This appendix provides complete formal specifications for each modality space.

\subsection{Graph Space $\Sigma_\mathsf{G}$}

A graph value $g \in \Sigma_\mathsf{G}$ is a labelled directed multigraph
$g = (N, E, \ell_N, \ell_E)$ where:
\begin{itemize}[leftmargin=*]
  \item $N \subseteq \mathsf{URI}$ is a set of node identifiers (IRIs in the
    RDF model, or opaque identifiers in the property-graph model).
  \item $E \subseteq N \times N$ is a set of directed edges.
  \item $\ell_N : N \to 2^{\mathsf{Label}}$ assigns zero or more labels to
    each node.
  \item $\ell_E : E \to \mathsf{Label} \times \mathsf{PropertyMap}$ assigns
    a type label and a property map to each edge.
\end{itemize}

For RDF compatibility, each edge $(s, o)$ with label $p$ corresponds to the
RDF triple $\langle s, p, o \rangle$.

\subsection{Vector Space $\Sigma_\mathsf{V}$}

A vector value $\vec{v} \in \Sigma_\mathsf{V}$ is a real-valued vector of
fixed dimension $d$:
\[
  \Sigma_\mathsf{V} = \mathbb{R}^d, \quad d \in \{128, 256, 384, 512, 768, 1024, 1536, 4096\}
\]
The dimension is fixed per octad store but may vary between federated stores.
Vectors are normalised to unit length when cosine similarity is the primary
distance metric.

\subsection{Tensor Space $\Sigma_\mathsf{T}$}

A tensor value $t \in \Sigma_\mathsf{T}$ is a typed multi-dimensional array:
\[
  \Sigma_\mathsf{T} = \bigcup_{n=1}^{N_\mathsf{max}}
    \bigcup_{\tau \in \mathsf{DType}}
    \tau^{d_1 \times d_2 \times \cdots \times d_n}
\]
where $\mathsf{DType} = \{\mathsf{f32}, \mathsf{f64}, \mathsf{i32}, \mathsf{i64},
\mathsf{bf16}\}$ and $N_\mathsf{max} = 8$ (maximum rank).

\subsection{Semantic Space $\Sigma_\mathsf{S}$}

A semantic value $s \in \Sigma_\mathsf{S}$ is a CBOR-encoded structure:
\[
  s = (\mathit{types} : \mathsf{List}[\mathsf{URI}],\;
       \mathit{proofs} : \mathsf{List}[\mathsf{Blob}],\;
       \mathit{policies} : \mathsf{List}[\mathsf{PolicyRef}])
\]
where $\mathsf{URI}$ references an external ontology, $\mathsf{Blob}$ is an
opaque byte sequence (typically a serialised proof term), and
$\mathsf{PolicyRef}$ is a hash-based reference to a policy octad.

\subsection{Document Space $\Sigma_\mathsf{D}$}

A document value $\mathit{doc} \in \Sigma_\mathsf{D}$ is a structured text
with inverted-index metadata:
\[
  \mathit{doc} = (\mathit{title} : \mathsf{String},\;
                  \mathit{body} : \mathsf{String},\;
                  \mathit{fields} : \mathsf{Map}[\mathsf{String}, \mathsf{String}])
\]
The inverted index is maintained incrementally; insertions and updates
trigger index segment merges according to Tantivy's merge policy.

\subsection{Temporal Space $\Sigma_\mathsf{P}$}

A temporal value $h \in \Sigma_\mathsf{P}$ is a Merkle tree:
\[
  h = \mathsf{MerkleTree}[\mathsf{Snapshot}]
\]
where each $\mathsf{Snapshot} = (\mathit{timestamp}, \mathit{state\_hash},
\mathit{parent\_hash})$ and $\mathit{state\_hash} =
\mathsf{SHA256}(\varphi(m_1) \| \cdots \| \varphi(m_8))$.

\subsection{Provenance Space $\Sigma_\mathsf{R}$}

A provenance value $c \in \Sigma_\mathsf{R}$ is a hash chain:
\[
  c = [e_1, e_2, \ldots, e_n]
\]
where $e_i = (\mathit{actor}, \mathit{action}, \mathit{timestamp},
\mathit{data\_hash}, \mathit{prev\_hash})$ and
$\mathit{prev\_hash} = \mathsf{SHA256}(e_{i-1})$.

\subsection{Spatial Space $\Sigma_\mathsf{X}$}

A spatial value $p \in \Sigma_\mathsf{X}$ is a WGS84 geometry:
\[
  \Sigma_\mathsf{X} = \mathsf{Point} \cup \mathsf{BBox} \cup \mathsf{Polygon}
\]
where $\mathsf{Point} = (\mathit{lat}, \mathit{lon}, \mathit{alt}?)$ with
$\mathit{lat} \in [-90, 90]$, $\mathit{lon} \in [-180, 180]$, and optional
altitude in metres.

\section{VQL Grammar (EBNF Excerpt)}\label{app:vql-grammar}

\begin{lstlisting}[language={},caption={VQL grammar excerpt (EBNF)},
  basicstyle=\ttfamily\scriptsize]
query       ::= SELECT projection FROM 'octads'
                WHERE condition+
                (ORDER BY ordering)?
                (LIMIT integer)?

projection  ::= proj_item (',' proj_item)*
proj_item   ::= 'octad' '.' modality '.' field
              | 'octad' '.' 'id'

condition   ::= vector_cond | graph_cond
              | spatial_cond | temporal_cond
              | document_cond | semantic_cond

vector_cond ::= 'octad' '.' 'vector' 'NEAR'
                '[' float (',' float)* ']'
                ('LIMIT' integer)?

graph_cond  ::= 'octad' '.' 'graph' 'MATCHES'
                '(' pattern ')'

spatial_cond::= 'octad' '.' 'spatial' 'WITHIN'
                spatial_fn

temporal_cond::= 'octad' '.' 'temporal' '.'
                 field  comp_op  value

document_cond::= 'octad' '.' 'document'
                  'CONTAINS' string_literal

spatial_fn  ::= 'radius' '(' lat ',' lon ','
                 distance unit ')'
              | 'bbox' '(' lat ',' lon ','
                 lat ',' lon ')'
\end{lstlisting}

\section{Drift Metric Derivations}\label{app:drift-derivations}

\subsection{Cosine Distance Sensitivity Analysis}

For the vector--document drift metric $d_{\mathsf{V},\mathsf{D}}$, we
analyse the sensitivity to document perturbations.  Let $\mathit{doc}$ be
the original document and $\mathit{doc}'$ a perturbed version with
$\|\mathit{doc}' - \mathit{doc}\|_\mathsf{edit} = k$ (edit distance).
For embedding functions based on transformer models, the relationship
between edit distance and cosine distance is approximately:
\[
  d_{\mathsf{V},\mathsf{D}} \approx
    1 - \exp\left(-\frac{k}{\lambda \cdot |\mathit{doc}|}\right)
\]
where $\lambda$ is an empirically determined sensitivity parameter
(typically $\lambda \approx 0.3$ for MiniLM-class models).  This
exponential relationship means that small edits produce small drift values,
while large rewrites produce drift values approaching 1, confirming the
metric's suitability as a staleness indicator.

\subsection{Jaccard Distance Monotonicity}

For the graph--semantic drift metric $d_{\mathsf{G},\mathsf{S}}$, we note
that the Jaccard distance is monotone in the symmetric difference:
\[
  |E_\mathsf{G} \triangle T_\mathsf{S}| \leq
  |E'_\mathsf{G} \triangle T_\mathsf{S}|
  \implies
  d_{\mathsf{G},\mathsf{S}}(g, s) \leq
  d_{\mathsf{G},\mathsf{S}}(g', s)
\]
This guarantees that adding edges to the graph without updating semantic
annotations strictly increases the drift value, providing a reliable
divergence signal.

\end{document}
